
This chapter connects to the classic roots of data analysis with low-dimensional structures: geometrically linear, statistically independent.

\section{Principal Component Analysis (PCA)} 
\subsection{Learning a linear subspace via denoising} 
Use the classic case of learning a single linear subspace to exemplify the general approaches to learn a low-dim model via compression (dimension reduction) from noising samples. Introduce the power-iteration algorithm... 

\subsection{Learning via completion and error correction}
Might worth mentioning variatnts such as Matrix Completion, Robust PCA for learning low-dimensional linear models under more general notion of ``noise'' or ``perturbantion''. 

\section{Independent Component Analysis (ICA)} Problem formulation... independence and incoherence, non-Gaussian... Kurtosis. Introduce the associated power-iteration algorithm.

\section{Sparse Coding and Dictionary Learning (DL)} Mixture of low-dim linear independent structures.... L1 and L4. 

\subsection{Sparse coding via ISTA} Introduce basic algorithm for sparse coding such as ISTA.

\subsection{Complete dictionary learning via $\ell^4$-maximization} 
our work on the complete case: \href{https://arxiv.org/abs/1906.02435}{L4-based dictionary learning} 

\subsection{Pursuit low-dimensionality via power iteration}
\href{https://openreview.net/pdf?id=SJeY-1BKDS}{The ICLR paper} provides a unifying view on PCA, ICA, and Dicionary Learning and their algorithms (all power-iteration type).

Interpretation of the algorithms as deep networks.




\section{Summary and Notes}
Unified computational mechanisms: Pursuit of low-dimensional linear and independent models via {\em power iteration} based on first-order oracles (gradients). 

Mention the work from John Wright and Qing Qu etc. on the over-complete dictionary learning.

Mention early generalizations to multiple subspaces (e.g. GPCA) and their limitations, and connections to later chapters on Rate Reduction.

